\documentclass[11pt,a4paper]{article}
\usepackage[T1]{fontenc}
\usepackage[utf8]{inputenc}
\usepackage{lmodern,microtype}
\usepackage[a4paper,margin=26mm,headheight=14pt]{geometry}
\usepackage{amsmath,amssymb,amsthm,mathtools}
\usepackage{booktabs,array,longtable}
\usepackage{xcolor,enumitem,listings,xurl,hyperref,fancyhdr,needspace}
\definecolor{ink}{HTML}{173047}
\definecolor{accent}{HTML}{176B70}
\definecolor{wash}{HTML}{F1F5F7}
\hypersetup{colorlinks=true,linkcolor=ink,citecolor=accent,urlcolor=accent,pdftitle={A reflexive-base introduction rule for iPRA},pdfauthor={Research report prepared for Vladimir Reshetnikov}}
\pagestyle{fancy}
\fancyhf{}
\fancyhead[L]{\small\textsf{iPRA / Reflexive-base introduction}}
\fancyhead[R]{\small\textsf{d9383325b18d}}
\fancyfoot[C]{\small\thepage}
\setlength{\parindent}{0pt}
\setlength{\parskip}{5pt}
\setlist{nosep,leftmargin=1.8em}
\setcounter{tocdepth}{2}
\newtheorem{theorem}{Theorem}[section]
\newtheorem{lemma}[theorem]{Lemma}
\newtheorem{proposition}[theorem]{Proposition}
\newtheorem{corollary}[theorem]{Corollary}
\theoremstyle{definition}
\newtheorem{definition}[theorem]{Definition}
\newcommand{\iPRA}{\mathrm{iPRA}}
\newcommand{\nf}{\operatorname{nf}}
\newcommand{\fnf}{\operatorname{fnf}}
\newcommand{\Teq}{T_{\mathrm{eq}}}
\newcommand{\file}[1]{\texttt{\detokenize{#1}}}
\newcommand{\Rcal}{\mathcal R}
\newcommand{\Zcal}{\mathcal Z}
\lstset{basicstyle=\ttfamily\scriptsize,columns=fullflexible,
  keepspaces=true,breaklines=true,showstringspaces=false,
  frame=single,framerule=.2pt,rulecolor=\color{black!20},
  backgroundcolor=\color{wash},xleftmargin=2pt,xrightmargin=2pt,
  aboveskip=8pt,belowskip=8pt}
\title{\color{ink}\textbf{A reflexive-base introduction rule for iPRA}\\[5pt]
  \large A six-rule presentation with ordinary closure semantics}
\author{Research report prepared for Vladimir Reshetnikov}
\date{11 November 2026}
\begin{document}
\maketitle
\thispagestyle{empty}

\begin{abstract}
Replace the current introduction rule's test on the syntactic form of its
existential witness by a test on its zero instance. Use equation progress
exactly when that instance is a syntactically reflexive equation; use numerical
cases otherwise. On the existing zero-normal term syntax this yields six
primitive rules and exactly the sequents of the repository's zero-normal
kernel. Missing case steps are recovered by cutting out a proved
reflexive base, and arbitrary equation induction is recovered through a
selector gate. The guard is preserved, but need not be reflected, by
substitution; the latter issue is handled explicitly. A raw-term version
works without normalization and has seven rules. Both versions inherit the
complete bi-interpretation with intuitionistic primitive recursive arithmetic,
with whole numerical domains, actual equality, and identity numerical
composites. The matched strict checker decreases from 3,095 to 3,084
normalized Python bytes and from 1,075 to 1,073 AST nodes; its 1,013-token
count is unchanged. This is a small, concrete simplification, not a global
minimality theorem or a large new checker-size record.
\end{abstract}

\section{Result and comparison boundary}

The comparison is pinned to ProveIt commit
\file{d9383325b18d349b8bfdc54ef564c5784cc029ec}, under
\path{Logic/PrimitiveRecursiveArithmetic/ConciseIPRA}.
The inspected README lists forty-nine research packages and the associated
Lean formalization~\cite{upstream}. At that snapshot the ordinary kernel
\file{NPrf} has eight rules, the witness-selected form has seven, and the
formally proved zero-normal kernel \file{NZPrf} has seven. The written
\file{zero_normal} package already combines zero normalization with the
witness-selected introduction to obtain six rules~\cite{zero}.

\textbf{The new proposal keeps that six-rule count.} Its change is the
introduction criterion
\[
  B(t,0)\text{ is syntactically reflexive},
\]
in place of the joint requirement that the matrix is an equation and the
witness is the literal zero. The term language, the meanings of the
logical binders, and the numerical domain stay the same. This criterion
directly identifies bases that can be proved before using any induction.

\begin{center}
\fcolorbox{accent}{wash}{\parbox{.9\linewidth}{
\textbf{Main theorem.}
For every normal context and formula,
\[
  \Gamma\vdash_{\Rcal} A
  \quad\Longleftrightarrow\quad
  \Gamma\vdash_{\Zcal} A.
\]
The proposed calculus is equiconsistent and bi-interpretable with full
first-order iPRA. The arithmetic interpretations have the whole domain
as their domain formula, actual equality, and numerical composite
graph \(x=y\).
}}
\end{center}

There are two separate simplicity claims. Mathematically, the rule no longer
requires a distinguished syntactic representative of the witness or the
construction of an alternative term denoting zero to recover a missing
case. Operationally, the base test reuses a formula the checker already
needs, and a validated-arity test gives a small matched source reduction.
The source reduction is only eleven bytes; calling it a major breakthrough
would overstate the result.

The construction does not beat the repository's separate 516-byte compact
guarded codec. That codec has a different certificate representation and
public syntax contract. The report does not conflate a fully validated
tagged checker with that smaller profile.

\subsection{What is new and what is inherited}

The zero-normal syntax, crossed one-term transport, numerical and logical
bridge to iPRA, and the general guarded-introduction strategy are inherited.
The specific reflexive-base test, its proof allowing a false guard to
become true under substitution, and its raw-term asymmetric selector gate
are the contribution of this report.

The cumulative rule retaining both the zero and current induction
assumptions is already present as \emph{Old DI} in
\file{inherited_choice_basis.tex}. It is not claimed as new here.
Neither deletion of the zero computation rule nor the six-rule count is
claimed as new. The comparison is against the pinned source, not
an unverified reconstruction of an earlier draft.

\subsection{The target theory}

Here iPRA means full intuitionistic first-order logic with equality, in a
language containing a symbol for each finite primitive-recursive
description: zero, successor, projections, composition, and primitive
recursion. Its nonlogical axioms are the corresponding defining equations,
successor separation and injectivity, and induction for quantifier-free
formulas, with parameters and retained assumptions. This is the target
defined in the repository's \file{PR.lean}; it is not merely the equational
fragment or a statement about standard natural-number evaluations.

\clearpage
\tableofcontents
\clearpage
